% !TEX TS-program = pdflatex
% !TEX encoding = UTF-8 Unicode

% This is a simple template for a LaTeX document using the "article" class.
% See "book", "report", "letter" for other types of document.

\documentclass[11pt, preview]{standalone} % use larger type; default would be 10pt

\usepackage[utf8]{inputenc} % set input encoding (not needed with XeLaTeX)


\usepackage{../../markup}
%%% Examples of Article customizations
% These packages are optional, depending whether you want the features they provide.
% See the LaTeX Companion or other references for full information.

%%% PAGE DIMENSIONS
\usepackage{geometry} % to change the page dimensions
\geometry{a4paper} % or letterpaper (US) or a5paper or....
% \geometry{margin=2in} % for example, change the margins to 2 inches all round
% \geometry{landscape} % set up the page for landscape
%   read geometry.pdf for detailed page layout information

\usepackage{graphicx} % support the \includegraphics command and options

% \usepackage[parfill]{parskip} % Activate to begin paragraphs with an empty line rather than an indent

%%% PACKAGES
\usepackage{amsmath, amsfonts,amssymb}
\usepackage{booktabs} % for much better looking tables
\usepackage{array} % for better arrays (eg matrices) in maths
\usepackage{paralist} % very flexible & customisable lists (eg. enumerate/itemize, etc.)
\usepackage{verbatim} % adds environment for commenting out blocks of text & for better verbatim
\usepackage{subfig} % make it possible to include more than one captioned figure/table in a single float
% These packages are all incorporated in the memoir class to one degree or another...

%%% HEADERS & FOOTERS
\usepackage{fancyhdr} % This should be set AFTER setting up the page geometry
\pagestyle{fancy} % options: empty , plain , fancy
\renewcommand{\headrulewidth}{0pt} % customise the layout...
\lhead{}\chead{}\rhead{}
\lfoot{}\cfoot{\thepage}\rfoot{}

%%% SECTION TITLE APPEARANCE
\usepackage{sectsty}
\allsectionsfont{\sffamily\mdseries\upshape} % (See the fntguide.pdf for font help)
% (This matches ConTeXt defaults)

%%% ToC (table of contents) APPEARANCE
\usepackage[nottoc,notlof,notlot]{tocbibind} % Put the bibliography in the ToC
\usepackage[titles,subfigure]{tocloft} % Alter the style of the Table of Contents

\renewcommand{\cftsecfont}{\rmfamily\mdseries\upshape}
\renewcommand{\cftsecpagefont}{\rmfamily\mdseries\upshape} % No bold!

\newcommand{\N}{\mathbb{N}}
\newcommand{\Z}{\mathbb{Z}}
\newcommand{\R}{\mathbb{R}}
\newcommand{\Q}{\mathbb{Q}}
%%% END Article customizations

%%% The "real" document content comes below...

%\date{} % Activate to display a given date or no date (if empty),
         % otherwise the current date is printed 

\begin{document}
\config{name}{Stable Matchings}
\noindent {\bf Stable Matchings}.

Say that we want to pair up the jobs (1,2,3) with the candidates (A,B,C); such a pairing is called a \emph{matching}. The table below gives the ranked preferences of the jobs and candidates, with most preferred on the left and least preferred on the right.

\begin{center}
\begin{tabular}{|c|c|}
\hline
job & ranking \\
\hline
1 & B, A, C \\
\hline
2 & A, B, C \\
\hline
3 & B, A, C \\
\hline
\end{tabular}
\qquad
\begin{tabular}{|c|c|}
\hline
candidate & ranking \\
\hline
A & 3, 2, 1\\
\hline
B & 1, 2, 3 \\
\hline
C & 2, 1, 3\\
\hline
\end{tabular}
\end{center}
Given a matching, we say that an unpaired job and candidate form a \emph{rogue couple} if both prefer each other to the partner they are currently matched to. 

\begin{enumerate}
\item Select all of the rogue couples for the matching (1A, 2B, 3C) with respect to the preferences in the table above.
\begin{Multi}
\Hint Review the definition of a matching and a rogue couple above. For more details read course note 4. Remember that a rogue couple requires the job and the candidate to {\it mutually} prefer each other!
\begin{enumerate}
\TrueChoice\item (1, B)
\FalseChoice\item (1, C)
\TrueChoice\item (2, A)
\FalseChoice\item (2, C)
\TrueChoice\item (3, A)
\FalseChoice\item (3, B)
\Solution Let's consider the choices in turn. 

a) Job $1$ ranks candidate $B$ above $1$'s current partner $A$, and candidate $B$ likewise ranks job $1$ above $B$'s current partner $2$, so $(1,\ B)$ {\it do} form a rogue couple.

b) Candidate $C$ prefers job $1$ over $C$'s current partner $3$, but job $1$ prefers their current partner $A$ over $C$, so $(1,\ C)$ do {\it not} form a rogue couple.

c) Job $2$ prefers $A$ over $2$'s current partner $B$, and candidate $A$ prefers $2$ over $A$'s current partner $1$, so $(2,\ A)$ {\it do} form a rogue couple.

d) Candidate $C$ prefers job $2$ over $C$'s current partner $3$, but job $2$ prefers their current partner $B$ over $C$, so $(2,\ C)$ do {\it not} form a rogue couple.

e) Job $3$ prefers $A$ over $3$'s current partner $C$, and candidate $A$ prefers $3$ over $A$'s current partner $1$, so $(3,\ A)$ {\it do} form a rogue couple.

f) Job $3$ prefers $B$ over $3$'s current partner $C$, but candidate $B$ prefers their current partner $2$ over $3$, so $(3,\ B)$ do {\it not} form a rogue couple.
\end{enumerate}
\end{Multi}
\item Again given the preferences above, select all of the rogue couples for the matching (3A, 2B, 1C).
\begin{Multi}
\Hint Review the definition of a matching and a rogue couple above. For more details read course note 4. Remember that a rogue couple requires the job and the candidate to {\it mutually} prefer each other!
\begin{enumerate}
\FalseChoice\item (1, A)
\TrueChoice\item (1, B)
\FalseChoice\item (2, A)
\FalseChoice\item (2, C)
\FalseChoice\item (3, B)
\FalseChoice\item (3, C)
\Solution Let's consider each choice in turn.

a) Job $1$ prefers candidate $A$ over $1$'s current partner $C$, but candidate $A$ prefers their current partner $3$ over $1$, so $(1,\ A)$ do {\it not} form a rogue couple.

b) Job $1$ and candidate $B$ both prefer each over their current partners, so $(1,\ B)$ {\it do} form a rogue couple.

c) Job $2$ prefers candidate $A$ over $2$'s current partner $B$, but candidate $A$ prefers their current partner $3$ over $2$, so $(2,\ A)$ do {\it not} form a rogue couple.

d) Candidate $C$ prefers job $2$ over $C$'s current partner $1$, but job $2$ prefers their current partner $B$ over $C$, so $(2,\ C)$ do {\it not} form a rogue couple.

e) Job $3$ prefers candidate $B$ over $3$'s current partner $A$, but candidate $B$ prefers their current partner $2$ over $3$, so $(3,\ B)$ do {\it not} form a rogue couple.

f) Job $3$ prefers their current partner $A$ over $C$, and candidate $C$ prefers their current partner $1$ over $3$, so $(3,\ C)$ do {\it not} form a rogue couple.
 
\end{enumerate}
\end{Multi}
%------------------------------------------------------------------------------------------------------
 A matching is considered \emph{stable} if it contains no rogue couples. This is because no one can convince anyone else to leave their partner in order to form a new pair. The propose-and-reject algorithm allows us to find stable matchings. This algorithm is described in course note 4. The following problems are designed to give you practice with this algorithm.
%-----------------------------------------------------
\item Consider the preference ranks given below for the jobs (1,2,3,4) and candidates (A,B,C,D). 
\begin{center}
\begin{tabular}{|c|c|}
\hline
job & ranking \\
\hline
1 & A, B, D, C \\
\hline
2 & A, C, B, D\\
\hline
3 & B, C, D, A\\
\hline
4 & B, A, D, C\\
\hline
\end{tabular}
\qquad
\begin{tabular}{|c|c|}
\hline
candidates & ranking \\
\hline
A & 3, 1, 4, 2\\
\hline
B &  2, 1, 3, 4\\
\hline
C & 4, 1, 3, 2\\
\hline
D & 1, 2, 3, 4\\
\hline
\end{tabular}
\end{center}
Simulate the propose-and-reject algorithm (with jobs proposing), and answer the following questions about the algorithm. We consider the first offers to have taken place on day 1.
\begin{enumerate}
\begin{Freeform}{A}
\item Which candidate did job $1$ propose to on day 1? %\ans A
% \Hint Review the stable marriage algorithm, found in course note 4. 
\Solution $1$'s top choice, $A$.
\end{Freeform}
\item Was job $1$ rejected? %\ans No
\begin{Choices}
\begin{itemize}
\FalseChoice \item Yes
\TrueChoice \item No
\Solution No, candidate $A$ received proposals from jobs $1$ and $2$, and $A$ preferred $1$ over $2$.
\end{itemize}
% \Hint Review the stable marriage algorithm, found in course note 4. 
\end{Choices}
\begin{Freeform}{B}
\item Which candidate did job $4$ propose to on day 1? %\ans B
% \Hint Review the stable marriage algorithm, found in course note 4.
\Solution $4$'s top choice, $B$.  
\end{Freeform}
\item Was job $4$ rejected on day 1? %\ans Yes
\begin{Choices}
\begin{itemize}
\TrueChoice \item Yes
\FalseChoice \item No
\Solution Yes, candidate $B$ received proposals from jobs $3$ and $4$, and $B$ preferred $3$ over $4$.
\end{itemize}
% \Hint Review the stable marriage algorithm, found in course note 4. 
\end{Choices}
\begin{Freeform}{4}
\item Which job proposed to candidate $A$ on day 2, that did not propose to $A$ on day 1? %\ans 4
% \Hint Review the stable marriage algorithm, found in course note 4. 
\Solution Job $4$, because $4$ was rejected by candidate $B$ on day 1, and $A$ is the next candidate on $4$ list.
\end{Freeform}
\item Was job $4$ rejected on day 2?
\begin{Choices}
\begin{itemize}
\TrueChoice \item Yes
\FalseChoice \item No
\Solution Yes, candidate $A$ received proposals from jobs $4$ and $1$, and $A$ preferred $1$ over $4$.
\end{itemize}
% \Hint Review the stable marriage algorithm, found in course note 4. 
\end{Choices}
\begin{Freeform}{A}
\item Which candidate was job $1$ matched to when the algorithm terminated? 
\Solution Candidate $A$.
\end{Freeform}
\begin{Freeform}{C}
\item Which candidate was job $2$ matched to when the algorithm terminated? %\ans B
% \Hint Review the stable marriage algorithm, found in course note 4. 
\Solution Candidate $C$.
\end{Freeform}
\begin{Freeform}{B}
\item Which candidate was job $3$ matched to when the algorithm terminated? %\ans D
% \Hint Review the stable marriage algorithm, found in course note 4.
\Solution Candidate $B$. 
\end{Freeform}
\begin{Freeform}{D}
\item Which candidate was job $4$ matched to when the algorithm terminated? %\ans A
% \Hint Review the stable marriage algorithm, found in course note 4. 
\Solution Candidate $D$.
\end{Freeform}

\item Now let us simulate the algorithm with candidates proposing. Before finding the outcome, what do you think is going to happen to the job's matches (compared to when the jobs were proposing)?
\begin{Choices}
\begin{itemize}
\FalseChoice\item Each job will either get the same match as before, or prefer their new match.
\TrueChoice\item Each job will either get the same match as before, or prefer their old match.
\FalseChoice\item Some jobs will get the same match as before, some will prefer their new matches, and some will prefer their old matches.
\end{itemize}
\Solution The algorithm run with jobs proposing produces a job-optimal pairing. Therefore the original matches are as high on their list as each job can hope for. Another way to arrive at the same conclusion is to note that the algorithm produces job-pessimal pairings if run with candidates proposing. Therefore jobs get partners that are as low as it can be on their list in the second run.
\end{Choices}

\item Which job will candidate $A$ propose to on the first day?
\begin{Freeform}{3}
\Solution $A$'s top choice, $3$.	
\end{Freeform}

\item Which candidate gets rejected on the first day?
\begin{Choices}
\begin{itemize}
\FalseChoice\item $A$
\FalseChoice\item $B$
\FalseChoice\item $C$
\FalseChoice\item $D$
\TrueChoice\item No one
\end{itemize}
\Solution No one gets rejected, because all jobs receive exactly one proposal.
\end{Choices}


\begin{Freeform}{D}
\item In a job-optimal stable matching, which candidate would be paired with job $4$?
\Solution The algorithm with jobs proposing returns a job-optimal pairing. We have already seen that the algorithm run with job proposing matches job $4$ with candidate $D$.
\end{Freeform}

\begin{Freeform}{C}
\item In a candidate-optimal stable matching, which candidate would be paired with job $4$?
\Solution The algorithm with candidates proposing returns a candidate-optimal pairing. In that matching, job $4$ is matched with candidate $C$.
\end{Freeform}

\begin{Freeform}{3}
\item In a job-pessimal stable matching, which job would be paired with candidate $A$?	
\Solution We get the job-pessimal pairing by running the algorithm with candidates proposing. In that case candidate $A$ is matched with job $3$.
\end{Freeform}

\begin{Freeform}{1}
\item In a candidate-pessimal stable matching, which job would be paired with candidate $A$?
\Solution We get the candidate-pessimal pairing by running the algorithm with jobs proposing. In that case candidate $A$ is matched with job $1$ as we have already seen.	
\end{Freeform}


\end{enumerate}
\end{enumerate}

\end{document}
